\chapter{新人如何安全做一次自己的回测实验}

\section{本章要解决的困惑}

读懂以后，下一步通常是“我想自己跑一次”。本章给的是安全实验路线，不是长 runbook。

原则：

\begin{lstlisting}
先复现已有 profile
再改一个小参数
再看 entries/exits/daily
最后才写研究结论
\end{lstlisting}

\section{实验前检查}

先确认你不是在修改 runtime 代码。第一轮只允许读数据、跑 diagnostic、看输出。

不要动：

\begin{lstlisting}
data/
date/
runs/   unless you are writing your own new run output
.env*
rpc.txt
\end{lstlisting}

不要把研究 label 接进 Bot 或 Runner。

\section{路线一：看 decision frame 是否存在}

目的不是训练模型，而是确认热路径存在：

\begin{lstlisting}
python systems/ccusdt_replay_exchange/diagnostics/decision_frame_cache.py validate ^
  --repo-root . ^
  --symbol CCUSDT ^
  --from-date 2026-05-18 ^
  --to-date 2026-05-18
\end{lstlisting}

它验证的是：

\begin{quote}
这一天的 \code{decision_frame_v1} cache schema、manifest、row count 和字段是否可用。
\end{quote}

\section{路线二：复现一个 fast profile}

跑 fast 时，明确写出：

\begin{lstlisting}
symbol
date range
exit profile
admission profile
capacity profile
leverage cap
threshold source
\end{lstlisting}

如果你只改一个参数，比如 \code{leverage-cap} 或 \code{gamma01}，输出目录名也要表达清楚。否则你几天后会不知道这次 run 到底是什么。

\section{路线三：看输出，而不是只看收益}

输出检查顺序：

\begin{enumerate}[leftmargin=2em]
  \item \code{summary.json}：profile 是否是你以为的；
  \item \code{daily.csv}：是不是一两天贡献了全部收益；
  \item \code{entries.parquet}：实际 entry 数和 cell 分布；
  \item \code{exits.parquet}：mid 和 executable 是否差很多；
  \item \code{profile_manifest.json}：admission/capacity/exit 是否记录完整。
\end{enumerate}

\section{路线四：只在必要时 strict}

如果你只是改了一个研究参数，先 fast。只有当你要验证运行路径时，才 strict：

\begin{lstlisting}
cargo run --manifest-path systems/ccusdt_replay_exchange/engine/Cargo.toml ^
  -p ccusdt_replay_cli -- run sparse-python ^
  --canonical-date 2026-05-18 ^
  --max-events 500 ^
  --latency-us 50000
\end{lstlisting}

这条命令验证的是本地交易所路径，不是大规模收益。

\section{实验记录模板}

每次实验至少写：

\begin{lstlisting}
Question:
  我想验证什么？

Runtime inputs:
  只读了哪些 runtime-safe 数据？

Profile:
  admission / capacity / exit / leverage / thresholds

Main readout:
  actual_entries, net_weighted_bps, daily_sign_rate, worst_day

Execution note:
  mid or top-of-book taker?

Conclusion:
  这只是研究证据，还是可进入下一轮 strict audit？
\end{lstlisting}

\section{常见误区}

\begin{itemize}[leftmargin=2em]
  \item 一次 run 好就说策略有效；
  \item 把 quote-only 窗口和 trade/L2 完整窗口混比；
  \item 只看 mid，不看 executable；
  \item 忘记 q70 是否 prior-day；
  \item 把研究报告里的 Delta/Energy 当成已上线逻辑。
\end{itemize}

\checkpoint{安全实验不是不探索，而是每次只改一个边界清楚的东西。}

